\chapter{Verification and Evidence Vocabulary}
\index{evidence!classes|textbf}
\index{verification|textbf}
\index{validation!scientific|textbf}

\label{ch:training-evidence-vocabulary}
\label{ch:verification-and-validation}

\chapterstatus{foundational evidence vocabulary. This chapter classifies claims
and records; it introduces no new numerical result, scientific validation, or
uncertainty-quantification result.}

\section{Learning goals}

After working through this chapter, a reader should be able to

\begin{enumerate}
  \item write a bounded claim before selecting evidence;
  \item distinguish mathematical proof, software verification, numerical
        verification, scientific validation, and uncertainty quantification;
  \item distinguish an oracle, metric, tolerance, observation, and disposition;
  \item distinguish a finite represented reference from an exact limit;
  \item separate transform training samples from withheld evaluation and
        independent reconstruction;
  \item keep parent-model, numerical, and model-reduction errors separate; and
  \item state what a passing check does not establish.
\end{enumerate}

Scientific software can run successfully, reproduce its own output, and still
answer the wrong mathematical or physical question. Evidence vocabulary is
therefore a computational prerequisite rather than an afterword.

\section{Claims precede evidence}
\index{claim!evidence requirement}

An evidence record begins with a bounded claim. The claim identifies the system,
quantity, representation, parameter domain, metric, tolerance, and intended use.
Only then can an oracle and acceptance rule be selected. A recurring record
shape is
\begin{equation}
  (\text{reference},\text{candidate},\text{metric},
   \text{tolerance},\text{observation},\text{disposition}).
\end{equation}
The record is incomplete unless the reference and candidate carry the metadata
required to make the metric meaningful.

These terms have distinct roles.

\begin{description}
  \item[Claim.] The bounded statement the evidence is intended to support.
  \item[Reference or oracle.] The independently justified source of expected
        behavior for the stated check.
  \item[Candidate.] The object or behavior under examination.
  \item[Metric.] The declared comparison quantity.
  \item[Tolerance or acceptance rule.] The prospectively stated decision rule.
  \item[Observation.] The value or categorical outcome actually obtained.
  \item[Disposition.] The conclusion permitted by the rule and evidence class.
\end{description}

A threshold-free diagnostic can report an observation without accepting or
rejecting it. Conversely, a tolerance without a named metric and scale is not a
complete acceptance rule.

\section{Evidence classes}

\begin{description}
  \item[Mathematical proof.] Establishes a theorem from explicitly stated
    assumptions by a recorded derivation or checked formal encoding. Proof does
    not establish that the assumptions describe the physical system, that a
    numerical implementation is correct, or that a bound is useful at project
    scale.
  \item[Software verification.] Establishes that software follows a documented
    contract. Examples include shape and type invariants, exact serialization
    round trips, basis-order checks, failure behavior, and declared symmetry or
    Hermiticity tests.
  \item[Numerical verification.] Establishes agreement between a numerical
    implementation and independently derived mathematics, or controlled
    convergence toward the stated mathematical problem. Examples include
    analytic limiting cases, manufactured solutions, independent norm
    calculations, and stencil-refinement studies.
  \item[Scientific validation.] Evaluates adequacy against independent trusted
    physical, experimental, or scientific reference evidence for a declared
    use. It addresses whether the selected model represents relevant physical
    behavior, not merely whether its code is internally consistent.
  \item[Uncertainty quantification.] Identifies relevant uncertainty sources and
    characterizes or propagates them to declared outputs. Sensitivity alone is
    not automatically a probabilistic uncertainty statement.
\end{description}

No class implies another. A proved implication does not validate its physical
assumptions. A passing software regression does not establish numerical
convergence. Agreement with a selected Kohn--Sham parent does not establish
agreement with experiment. A validation comparison does not quantify every
uncertainty source.

\citationtodo{high priority}{Check the evidence-class taxonomy and the claim
that the classes do not imply one another against an authoritative
verification, validation, and uncertainty-quantification source. Candidate:
National Research Council (2012), provisional key
\texttt{nationalResearchCouncil2012}; the candidate record is present for this
editorial note.}

\section{Examples that share an array but not a claim}

A Hermitian matrix may participate in several different evidence records.

\begin{itemize}
  \item Comparing \(H\) with \(H^{\dagger}\) can verify a declared software
        invariant.
  \item Comparing \(H\) with an independently derived small-matrix oracle can
        numerically verify an implementation.
  \item Comparing an observable derived from \(H\) with independent physical
        data can contribute to scientific validation for a declared use.
\end{itemize}

The array does not carry an evidence class intrinsically. The class follows
from the claim, reference, construction, and interpretation.

Similarly, successful process completion is an observation about execution. It
is not automatically evidence that inputs were scientifically accepted, a
solver converged to the desired state, a discretization converged, or the model
is valid.

\section{Errors from different operations}
\index{error!category separation}

Projection error, discretization error, fitting error, truncation error, and
parent-model error arise from different operations. They should be combined
only when an owning mathematical contract defines a compatible relationship.
The broad categories used throughout this monograph are

\begin{description}
  \item[Parent-model error.] Discrepancy attributable to the selected physical
        parent and its assumptions relative to the intended physical use.
  \item[Numerical or discretization error.] Discrepancy introduced while
        approximating the stated mathematical problem at finite settings.
  \item[Model-reduction error.] Discrepancy introduced when replacing the
        accepted parent representation by a reduced class.
\end{description}

A reduced model can reproduce an unconverged parent accurately. A converged
numerical solution can solve an inadequate parent model precisely. Agreement
with experiment can arise through cancellation between parent and reduction
errors. Evidence records must preserve these possibilities rather than report
one undifferentiated ``model error.''

\section{Finite-setting observations}
\index{convergence!finite-setting}

Suppose a quantity \(q_j\) is computed at two successive finite settings and
satisfies
\begin{equation}
  |q_{j+1}-q_j|\leq\tau.
\end{equation}
This establishes stability over those tested settings under the recorded
controls. It does not by itself bound
\begin{equation}
  |q_j-q_{\infty}|.
\end{equation}
An infinite-setting bound requires additional structure, such as a justified
asymptotic error model, monotonicity result, variational bound, or sufficient
higher-resolution evidence. ``Stable over the tested sequence'' is therefore a
stronger description than ``converged'' when only the finite contrast is known.

Derived quantities can amplify errors. A central second-difference estimate of
curvature divides an energy combination by a squared step. Pointwise energy
agreement therefore need not control effective mass. Conditioning and
cancellation must be evaluated for the derived claim, not inferred from a
visually smooth parent plot.

\section{Finite references and exact language}
\index{reference!finite represented}

A numerical study may compare candidates with a declared high-resolution
finite representation. If $q_h$ is the candidate and $q_{h_{\mathrm{ref}}}$ is
that reference, the observed quantity
\begin{equation}
  \epsilon_{h,\mathrm{ref}}
  =\left|q_h-q_{h_{\mathrm{ref}}}\right|
\end{equation}
is a finite-reference discrepancy. It is not automatically an error bound to
the continuum quantity $q$. The latter would require information about
$|q_{h_{\mathrm{ref}}}-q|$ or a theorem that otherwise controls the limit.

This distinction changes the permitted language. A decreasing sequence of
finite-reference discrepancies can be described as observed refinement over
the tested settings. Calling the final value an exact continuum error would
state more than the comparison measured. Small nonmonotone variations near
floating-point or iterative-solver scales should also be retained rather than
removed to manufacture a monotone convergence plot.

\section{Transform training data and withheld evaluation}
\index{training data!finite transform}
\index{withheld evaluation!role separation}

Training data can play a mathematical role beyond statistical parameter
estimation. In a complete finite Fourier transform, the training mesh defines
all transform coefficients. Exact reconstruction on that same mesh verifies
the forward--inverse transform pair, but it is interpolation of the defining
samples. It does not test behavior between them.

A withheld mesh can evaluate the frozen interpolant at coordinates that were
not used to calculate coefficients. To preserve that role, the withheld values
must not alter
\begin{itemize}
  \item the fitted coefficients;
  \item the candidate model class or hopping range;
  \item the stopping rule;
  \item the metric or tolerance; or
  \item the selection of which outcomes are reported.
\end{itemize}
If withheld observations motivate redesign, the original observation remains
part of the evidence history and a new untouched evaluation design is needed
for a later confirmatory claim. In a synthetic numerical study, disjoint
withheld evaluation remains numerical verification relative to the selected
parent. The word ``withheld'' does not promote it to scientific validation.

\section{Independent reconstruction is not producer replay}
\index{verification!independent reconstruction}

A producer may verify internal invariants before returning a result. That is
valuable, but importing the producer and asking it to generate the same object
again can repeat the same mistake. Stronger numerical-consistency evidence
reconstructs decisive quantities from frozen inputs and retained outputs by a
separately maintained route.

Independence is bounded rather than absolute. A verifier may share the same
language runtime, numerical libraries, or mathematical assumptions. Its record
should therefore state what it reconstructs independently and what remains
shared. A useful verifier also checks identities and roles---for example,
coordinate inventories, transform representatives, tolerances, and the
separation of training and withheld samples---rather than comparing only one
headline scalar.

Agreement within a prospectively frozen numerical tolerance establishes
consistency for the reconstructed finite protocol. It does not select a model
for a scientific use, quantify physical uncertainty, or validate the parent
model. Chapter~\ref{ch:training-isolated-band-fourier-reduction} supplies a
concrete transform, truncation, and route-comparison setting to which these
rules apply.

\section{Negative evidence and stop conditions}
\index{negative evidence}
\index{stop condition}

A failed acceptance criterion is retained rather than weakened. It may indicate
an implementation defect, inadequate resolution, incompatible representation,
insufficient model-class expressiveness, or mismatch between a model and its
intended use. These diagnoses have different remedies.

A missing prerequisite can also be a valid stop condition. If units, geometry,
basis ordering, energy reference, or alignment are incompatible, direct matrix
subtraction should not proceed to a residual. A fail-closed compatibility
finding is evidence about the attempted operation, not an inconvenience to be
removed by guessing metadata.

A well-executed negative result is bounded by the tested domain. Failure of one
frozen model hierarchy does not prove failure of every possible model, and one
nonconverged numerical sequence does not prove that convergence is impossible.

\section{Reproducibility does not promote evidence}
\index{reproducibility}
\index{provenance}

A result is reproducible only relative to retained identities and declared
conditions. A useful record generally includes immutable inputs, code and
dependency versions, physical and numerical settings, commands, artifact
checksums, metric definitions, structured outcomes, warnings, and locations of
externally retained large artifacts.

Reproduction preserves the class of the reproduced claim. Reproducing a
pedagogical calculation reproduces pedagogical behavior; reproducing a
synthetic verification case reproduces bounded verification evidence;
reproducing a production calculation reproduces that calculation under its
recorded conditions. None is promoted automatically to scientific validation.

\section{A compact evidence ledger}

For each computational stage, a ledger can record

\begin{enumerate}
  \item the claim identifier and bounded statement;
  \item identities of reference and candidate records;
  \item compatibility checks and stop conditions;
  \item metric, unit, normalization, and tolerance;
  \item observed value and pass, fail, or diagnostic-only disposition;
  \item evidence class and limitations; and
  \item provenance for reproduction.
\end{enumerate}

This structure makes limitations part of the result rather than prose added
only after a surprising outcome.

\section{Exercises}

\begin{enumerate}
  \item Classify a JSON round trip, a manufactured finite-difference solution,
        and comparison with an experimental band edge by evidence class.
  \item Write a bounded claim for a Hermiticity check, including metric, unit,
        and tolerance. State three conclusions it does not support.
  \item Explain why two stable successive settings do not establish a bound to
        an infinite-setting limit.
  \item Give one example of parent-model, numerical, and model-reduction error
        in a single reduction workflow.
  \item Design a negative-evidence record for incompatible energy references
        that stops before matrix subtraction.
\end{enumerate}

\section{Next use and evidence boundary}

Part~II applies this vocabulary to model adequacy, aligned operator comparison,
and the retained benchmark inventory. The vocabulary itself does not establish
that any project calculation has passed. Evidence status remains owned by the
applicable proof package, test record, calculation record, scientific review,
or accepted disposition.
